%-*-tex-*-
\ifundefined{writestatus} \input status \relax \fi %
\chcode{docsty}

\def\cqu{}

\chapterhead{docsty}{DOCUMENT\cr STYLES}
\intex\ supplies a number of general and specialized document styles. It is
also easy to create others within the basic structure given. 
The two basic styles supported by \intex\ is the (default) |\paperstyle| and
the |\bookstyle| of which this is an example. Different styles are usually
distinguished by different header/footer conventions, appendix placement and
numbering, and perhaps default fonts. 

In addition to the page changes, a few special commands exist that are
introduced  in the various styles. Since the |\paperstyle| the \intex\
default, all the |\title|, and abstract commands defined for that are
available, and will work, for {\bf all} styles.

\shead{paperstyle}{Paperstyle}
|\paperstyle| is the default style for \intex\. There will be no messages
indicating that you are using this style. This style is designed for 8.5in 
by 11in paper, |\tenpoint| fonts, no header, and a footer that consists of
the page number centered at the bottom in this style ``- 6 -'', where 6 would
have been the page number. In addition to the normal features mentioned
previously, there is a |\title| command and a 
|\beginabstract...\endabstract| pair. 

The description of the commands are:

\beginblockmode
\ext\@|\beginabstract...\endabstract|
\nbr
This command should enclose the entire abstract. The abstract font is the
|\smaller| font and both sides are indented. The text is automatically
prefaced by ``Abstract:'' in English and ``Sommaire:'' in French. 

\mbr
\ext\@|\title{<first row>\cr ... <last row>}|
\nbr
If the total length of the title is more than one line it is automatically
broken into multiple lines with each line centered. However, if |\cr| are
included as part of the title text, the lines will be broken there and
centered. The title font, called |\titlefont| is  |\bigfont|.
\endblockmode

\shead{inrsstyle}{INRS Report Style}
This style is very similar to the basic |\paperstyle|. The only addition is a
set of commands that create a cover. This style is called with either
\@|\formatinrs| or \@|\inrsreportstyle|. The title, report date, author list,
and report number are token lists. The actual cover is made with the command
|\faitepagetitre|.  An example of complete cover page is
\begintt
\formatinrs
\titrederapport = {INRS\TeX REFERENCE BOOK}
\numeroderapport = {84-19}
\datederapport = {juin 1984}
\auteursderapport = {Michael J. Ferguson}
\faitepagetitre
\ejectpage
\done
\endtt

\ejectpage

\done
